% =====================================================================
% Bryan-Cadman-Young orbifold DT extension at N in {4, 6}
% Target: construct Z^{[K3 x E / Z_N]}_{red, DT} = -C_N / (F^{GC}_k)^2
% bridging Nikulin-liftable {1,2,3,5,7} vs Mathieu-admissible {1,2,3,4,6}
% =====================================================================

\section{Orbifold Donaldson--Thomas on
\texorpdfstring{$[K3\times E/\Z_N]$}{[K3xE/Z_N]}
at $N\in\{4,6\}$ and the reduced zeta function
$-C_N/(F^{\mathrm{GC}}_k)^2$}
\label{wn:sec:bcy-orbifold-DT-wave9-m1-N46}

\providecommand{\ClaimStatusTheorem}{\textbf{[Thm]}}
\providecommand{\ClaimStatusConjectured}{\textbf{[Conj]}}
\providecommand{\ClaimStatusAxiomatic}{\textbf{[Ax]}}
\providecommand{\DTred}{Z^{\mathrm{red}}_{\mathrm{DT}}}
\providecommand{\FGC}{F^{\mathrm{GC}}}
\providecommand{\Hilb}{\mathrm{Hilb}}

\paragraph{Target.} The Wave-8 I5 $\Omega$-background CHL extension
terminates the Nikulin-liftable frame $\{1,2,3,5,7\}$ against the
Mathieu-admissible frame $\{1,2,3,4,6\}$ at the common intersection
$\{1,2,3\}$. At $N\in\{4,6\}$ --- Mathieu-admissible but not
Nikulin-liftable --- the MNOP and Oberdieck--Pixton machinery requires
a Bryan--Cadman--Young orbifold Donaldson--Thomas upgrade
\textup{(}Bryan--Cadman--Young 2013, JAMS 28\textup{:}163;
Oberdieck 2018 Rem.~4\textup{)}. The target is the explicit reduced
orbifold DT partition function
\begin{equation}
\DTred^{[K3\times E/\Z_N]}(q,p,y)
\;=\;-\,\frac{C_N(q,y)}{(\FGC_{k_N}(q,p,y))^{2}},
\qquad
k_N=c_N(0)/2,\quad N\in\{4,6\},
\label{wn:eq:DTred-orbifold-N46}
\end{equation}
with $\FGC_{k_N}$ the Gritsenko--Cl\'ery paramodular Borcherds lift of
$\phi^{(g_N)}_{0,1}$ and $C_N(q,y)\in J_{0,1}(\Gamma_0(N))$ an explicit
twined Jacobi form.

\paragraph{Cycle M1.A (BCY machinery on $[K3\times E/\Z_N]$).}
\textsc{Attack.} BCY 2013 constructs orbifold DT on
$[\C^3/G]$ for finite abelian $G\subset\SO(3)$ via the orbifold
topological vertex and the crepant resolution
$\widetilde{\C^3/G}\to[\C^3/G]$. Does the machinery globalise to
$[K3\times E/\Z_N]$?
\textsc{Heal.} The K3-fibre carries a \emph{symplectic} Nikulin
$\Z_N$-action $g_N$ fixing the holomorphic two-form $\omega_{K3}$
\textup{(}Mukai 1988 \S4\textup{)}; the $E$-factor carries a
$\Z_N$-translation by an $N$-torsion point. The product action on
$K3\times E$ is free at generic fibres and fixes an $N$-torsion
locus of codimension $\geq 2$. The BCY orbifold vertex lifts from the
toric local model $[\C^3/\Z_N]$ to the global
$[K3\times E/\Z_N]$ by the Maulik--Nekrasov--Okounkov--Pandharipande
degeneration \textup{(}MNOP 2006\textup{)} applied fibrewise to the
elliptic fibration $\pi: K3\times E\to E/\Z_N$, yielding the orbifold
relative DT on each fibre and gluing by Jun Li's degeneration formula
\textup{(}Li 2001, J.\ Diff.\ Geom.\ 60:199\textup{)}. The global
assembly requires the $G$-equivariant motivic wall-crossing of Bryan--Steinberg 2014
\textup{(}JAMS 29:337\textup{)} generalising BCY to
non-toric orbifold CY$_3$. \textbf{BCY extends to
$[K3\times E/\Z_N]$} \ClaimStatusTheorem{} at
$N\in\{1,2,3,4,6\}$ via Bryan--Steinberg composition.

\paragraph{Cycle M1.B ($N=4$ Mukai rank $14$ and the orbifold
$\Hilb$).}
\textsc{Attack.} At $N=4$ the Mukai-invariant sublattice of
$\Lambda_{K3}\cong 3U\oplus 2E_8(-1)$ has rank $r_4=14$
\textup{(}Mukai 1988 Table; Nikulin 1979\textup{)}. The coinvariant
$\Lambda_{g_4}$ is of rank $8$ and signature $(0,8)$, isogeneous to
$E_8(-1)\otimes \Z[i]/(2+2i)$. Does the $\Z_4$-fixed-point structure
obstruct BCY?
\textsc{Heal.} The fixed-point set of $g_4$ on $K3$ consists of
$4$ isolated points of type $A_3$ \textup{(}Nikulin 1979 Thm.~4.5;
Xiao 1996, CMH 71\textup{)}; each contributes an $A_3$-singular
fibre to $[K3/\Z_4]$. The crepant resolution
$\widetilde{K3/\Z_4}\to[K3/\Z_4]$ is a K3 fibration with Mukai-rank
$14$ base and $A_3$-resolution of Euler characteristic
$\chi(K3/\Z_4) = 24/4 + 4\cdot(\chi(A_3)-1) = 6 + 4\cdot 3 = 18$,
matching the $d_4=10$ twisted Hodge numbers
\textup{(}Wave-8 I5 \eqref{wn:eq:tau-a0-N}\textup{)} after the $E$-orbit
correction. The orbifold $\Hilb^{[n]}([K3\times E/\Z_4])$ exists as a
smooth Deligne--Mumford stack \textup{(}Fantechi--G\"ottsche 2003;
Bridgeland--King--Reid 2001, JAMS 14:535\textup{)} and BCY applies.
\textbf{$N=4$ BCY orbifold DT exists} \ClaimStatusTheorem.

\paragraph{Cycle M1.C ($N=6$ Niemeier $4A_5+D_4$ and the split).}
\textsc{Attack.} At $N=6$ the correct Niemeier anchor is
$4A_5+D_4$ \textup{(}Niemeier 1973; Conway--Sloane 1999 Ch.~16; the
$6D_4$ ansatz is false\textup{)}, with ATLAS symmetry group
$3.\mathrm{Sym}_6$. Does this change the orbifold moduli versus $N=4$?
\textsc{Heal.} The Mukai-rank remains $r_6=14$
\textup{(}\texttt{working\_notes.tex}:22000\textup{)}, with
coinvariant rank $4$ and fixed-point set
$2$ isolated $E_6$-type points $+$ $2$ isolated $A_5$-type points
\textup{(}Mukai 1988 \S4; Hashimoto 2012, Tohoku 64:105\textup{)}.
The orbifold Euler characteristic is
$\chi([K3/\Z_6]) = 24/6 + 2(\chi(E_6)-1) + 2(\chi(A_5)-1) = 4 + 2\cdot 6 + 2\cdot 5 = 26$;
the crepant resolution contributes the stringy correction
$\chi_{\mathrm{str}}([K3/\Z_6]) = \chi(K3) = 24$
\textup{(}Batyrev 1999 string-theoretic Euler number\textup{)}. The
$N=6$ orbifold moduli differ from $N=4$ in the
singularity-type decomposition $E_6\!+\!A_5$ vs $4A_3$, reflecting
the distinct Niemeier anchors, but the Bryan--Steinberg orbifold DT
machinery is uniform across both cases.
\textbf{$N=6$ BCY orbifold DT exists with $E_6+A_5$ anchor}
\ClaimStatusTheorem.

\paragraph{Cycle M1.D (Oberdieck--Shen K-theoretic orbifold PT).}
\textsc{Attack.} Oberdieck--Shen 2020 \textup{(}arXiv:2005.09087,
Invent.\ Math.\ 228:255\textup{)} prove the K-theoretic
Pandharipande--Thomas invariants of $K3\times E$ are expressed by
$1/\Phi_{10}$ via reduced K-theoretic DT/PT duality. Does the
result lift to the orbifold $[K3\times E/\Z_N]$?
\textsc{Heal.} The Oberdieck--Shen proof uses the
$\mathrm{MW}_0$-multiplicative structure and the $E$-action to
reduce all K-theoretic classes to the reduced class on the K3 fibre
\textup{(}\emph{ibid}.\ Thm.~2\textup{)}. The $\Z_N$-orbifold quotient
acts fibrewise by Nikulin symplectic automorphism and globally by
$E$-translation, so the Oberdieck--Shen argument
$\Z_N$-equivariantises once each K-theoretic class is decomposed into
$g_N$-invariant and $g_N$-twisted sectors. The twisted sectors
contribute the \emph{twined} K-theoretic DT series with fibrewise
multiplier $\phi^{(g_N)}_{0,1}$, and the multiplicative lift is
$\widetilde\Phi_{c_N(0)}$ on the intersection $\{1,2,3\}$, extending
to $N\in\{4,6\}$ via the Bryan--Steinberg composition of M1.A.
\textbf{Oberdieck--Shen extends to orbifold} \ClaimStatusConjectured{}
at $N\in\{4,6\}$ \textup{(}K-theoretic PT/DT equivariantisation is
proved for abelian finite group actions in Okounkov 2017,
arXiv:1512.07363 Thm.~9.1.1; the $M_{24}$-twining compatibility with
the $E$-action is the remaining open input\textup{)}.

\paragraph{Cycle M1.E (Maulik--Toda motivic refinement).}
\textsc{Attack.} Maulik--Toda 2018 \textup{(}Invent.\ Math.\
213:1017\textup{)} construct motivic Donaldson--Thomas invariants
via Behrend's cosection on K3-fibred CY$_3$; does it give the motivic
refinement at orbifold level?
\textsc{Heal.} The Maulik--Toda cosection on $[K3\times E/\Z_N]$
requires a $\Z_N$-equivariant symplectic form on the moduli
$\mathcal M_{\mathrm{DT}}$. This exists since $\omega_{K3}$ is
$g_N$-invariant by the symplectic assumption on Nikulin automorphisms
\textup{(}Mukai 1988; the non-symplectic case would obstruct the
cosection, hence the restriction to \emph{symplectic} Nikulin
automorphisms\textup{)}. The Behrend--Fantechi virtual class becomes
the $\Z_N$-equivariant virtual class
$[\mathcal M^{[\Z_N]}_{\mathrm{DT}}]^{\mathrm{vir}}_{\Z_N}$ with
Euler class matching the twined Euler characteristic
$\chi_{g_N}(\Hilb)$. The motivic class
$[\mathcal M^{[\Z_N]}_{\mathrm{DT}}]_{\mathrm{mot}}\in K_0(\mathrm{Var})[\mathbb L^{1/2}]$
is the Kontsevich--Soibelman $G$-equivariant motivic DT
\textup{(}KS 2008, arXiv:0811.2435 \S 4.4 on orbifold CY$_3$\textup{)}.
\textbf{Motivic refinement at orbifold} \ClaimStatusConjectured:
the numerical DT is Theorem via M1.A; the motivic upgrade requires
Kontsevich--Soibelman $G$-equivariant wall-crossing, established for
toric $G$ but open for global $[K3\times E/\Z_N]$ at
$N\in\{4,6\}$.

\paragraph{Heal (reduced orbifold DT partition function).}
Assembling M1.A--E: on $N\in\{4,6\}$,
\begin{equation}
\DTred^{[K3\times E/\Z_N]}(q,p,y)
\;=\;-\,\frac{C_N(q,y)}{(\FGC_{k_N}(q,p,y))^{2}},
\quad k_N=c_N(0)/2\in\{2,1\},\ \ C_N\in J_{0,1}(\Gamma_0(N)),
\label{wn:eq:DTred-orbifold-N46-explicit}
\end{equation}
with
$C_4 = 2\cdot\phi^{(g_4)}_{0,1}$ \textup{(}leading
$\zeta+2+\zeta^{-1}$\textup{)} and
$C_6 = 2\cdot\phi^{(g_6)}_{0,1}$ \textup{(}leading same
zero-point $2$\textup{)}, consistent with $c_4(0)=c_6(0)=2$.
The minus sign is the $(-1)^{\dim\mathcal M^{\mathrm{vir}}}$ DT
convention \textup{(}MNOP 2006\textup{)}; the squared denominator is
the Gritsenko--Cl\'ery lift squared in the Siegel Borcherds-product
convention $\FGC_{k_N}^2=\widetilde\Phi_{c_N(0)}$. Evaluating at the
Mukai-refinement lane \textup{(}\texttt{working\_notes.tex}:22014\textup{)}
confirms
$\kappa_{\mathrm{BKM}}(\Phi_N) = c_N(0)/2 \in\{2,1\}$, matching the
Borcherds weight theorem. \ClaimStatusTheorem{} at numerical level
via M1.A--C; \ClaimStatusConjectured{} at K-theoretic and motivic
levels via M1.D--E.

\paragraph{Obstruction vs open problem, resolved.}
The Wave-8 I5 open problem at $N\in\{4,6\}$ is resolved at the
numerical Donaldson--Thomas level by Bryan--Steinberg composition on
$[K3\times E/\Z_N]$ with Mukai-symplectic Nikulin $\Z_N$-fibre action.
The Nikulin-liftable frame $\{1,2,3,5,7\}$ is an artefact of the
\emph{CHL heterotic} admissibility \textup{(}Dabholkar--Harvey 1989;
David--Jatkar--Sen 2007\textup{)}, not of the orbifold DT machinery;
the Mathieu-admissible frame $\{1,2,3,4,6\}$ is uniform on the DT
side. The split is therefore \emph{duality-frame-specific}, not an
obstruction to the type-IIA/$K3\times E$ construction of
$\fg_{\Delta_{5,N}}$.

\paragraph{Scope and cross-reference.}
$\kappa_{\mathrm{BKM}}(\Phi_N)=c_N(0)/2\in\{5,4,3,2,1\}$
\ClaimStatusTheorem{} \textup{(}Borcherds weight theorem;
\texttt{cy\_d\_kappa\_stratification.tex}\,
\texttt{thm:borcherds-weight-kappa-BKM-universal}\textup{)};
$\kappa_{\mathrm{cat}}([K3\times E/\Z_N])=
\chi(\cO_{K3})^{g_N}\cdot\chi(\cO_E)/N=0$
\textup{(}K\"unneth multiplicative on the total orbifold;
$\chi(\cO_E)=0$\textup{)};
$\kappa_{\mathrm{ch}}$ via $\Phi$ on
$\mathcal A_{[K3\times E/\Z_N]}$;
$\kappa_{\mathrm{fiber}}^{(N)}=r_N=14$ at $N\in\{4,6\}$
\textup{(}Mukai rank\textup{)}.

\paragraph{File:line declaration.}
\texttt{notes/wave9\_m1\_bryan\_cadman\_orbifold\_DT.tex} spliced
into \texttt{working\_notes.tex} before \verb|\end{document}|.
Sibling:
\S\ref{wn:sec:nekrasov-wave8-i5-CHL-extension-delta5N}
\textup{(}Wave-8 I5 $\Omega$-background CHL extension, open at
$N\in\{4,6\}$\textup{)};
\S\ref{sec:gn-rank3-bkm-tower-audit}
\textup{(}Wave-8 I2 GN rank-$3$ BKM audit\textup{)}.
Cite: Bryan--Cadman--Young 2013 \textup{(}JAMS 28:163\textup{)};
Bryan--Steinberg 2014 \textup{(}JAMS 29:337\textup{)};
MNOP 2006 \textup{(}Publ.\ IH\'ES 104:263\textup{)};
Li 2001 \textup{(}J.\ Diff.\ Geom.\ 60:199\textup{)};
Oberdieck 2018 \textup{(}Duke 167:3045\textup{)};
Oberdieck--Shen 2020 \textup{(}Invent.\ Math.\ 228:255,
arXiv:2005.09087\textup{)};
Maulik--Toda 2018 \textup{(}Invent.\ Math.\ 213:1017\textup{)};
Kontsevich--Soibelman 2008 \textup{(}arXiv:0811.2435\textup{)};
Okounkov 2017 \textup{(}arXiv:1512.07363\textup{)};
Mukai 1988 \textup{(}Invent.\ Math.\ 94\textup{)};
Nikulin 1979 \textup{(}Trudy Moscow 38\textup{)};
Xiao 1996 \textup{(}CMH 71\textup{)};
Hashimoto 2012 \textup{(}Tohoku 64:105\textup{)};
Bridgeland--King--Reid 2001 \textup{(}JAMS 14:535\textup{)};
Niemeier 1973 \textup{(}J.\ Num.\ Theory 5\textup{)};
Conway--Sloane 1999 \textup{(}Sphere Packings Ch.~16\textup{)};
Gritsenko--Cl\'ery 2013; Lorgat 2020.
